\section{Thiết kế cuối cùng}

\subsection{Bản mẫu hoàn thiện}
Tiếp thu các ý kiến đóng góp quý báu từ đợt kiểm thử khả năng sử dụng ở Chương 7, nhóm đã thực hiện tinh chỉnh và hoàn thiện phiên bản Thiết kế cuối cùng (Final Refined Design):
\begin{itemize}
    \item \textbf{Tích hợp ô ghi chú dặn dò y tế trực tiếp vào luồng Đặt lịch}: Thay vì bắt người dùng phải điều hướng vào trang Hồ sơ riêng, hệ thống đưa trực tiếp trường nhập \textit{''Lưu ý y tế \& Dị ứng da''} vào Bước 3 của Stepper đặt lịch, đồng thời tự động điền dữ liệu đã lưu từ hồ sơ thú cưng để người dùng chỉ cần xác nhận hoặc chỉnh sửa nhanh.
    \item \textbf{Tăng cường độ tương phản cho Ma trận khung giờ}: Điều chỉnh màu sắc cho các slot đã kín chỗ sang màu xám đậm có dấu gạch chéo mờ và nhãn \textit{''Hết chỗ''}, các slot còn trống được viền xanh dương đậm nổi bật, giúp giảm thiểu tối đa sự nhầm lẫn thị giác.
    \item \textbf{Thêm biểu tượng phân biệt trực quan cho các mốc tiến độ}: Bổ sung các biểu tượng động rõ nét cho từng mốc (Mốc 1: Tiếp nhận -- Biểu tượng Bàn tay/QR; Mốc 2: Đang chăm sóc -- Biểu tượng Bọt xà phòng/Kéo cắt tỉa nhấp nháy; Mốc 3: Hoàn tất -- Biểu tượng Dấu tích xanh; Mốc 4: Chờ đón -- Biểu tượng Chuông thông báo), giúp người dùng nhận biết ngay lập tức chỉ trong $1$ giây nhìn lướt.
    \item \textbf{Tự động sắp xếp Lịch sử theo thứ tự mới nhất}: Đặt mặc định hiển thị phiên chăm sóc gần nhất ở đầu danh sách và bổ sung bộ lọc theo tháng/dịch vụ.
\end{itemize}

\subsection{Các cải tiến thiết kế then chốt}
Bảng tổng hợp các cải tiến then chốt của thiết kế hoàn thiện so với quy trình truyền thống:
\begin{table}[H]
\centering
\begin{tabularx}{\textwidth}{|X|X|X|}
\hline
\textbf{Vấn đề quy trình cũ (Proposal)} & \textbf{Giải pháp thiết kế mới} & \textbf{Cải thiện trải nghiệm đạt được} \\ \hline
Chờ xác nhận lịch hẹn lâu, dễ bị trôi tin nhắn hoặc kín chỗ. & Đặt lịch trực tuyến qua ma trận giờ trống thời gian thực \& Xác nhận tức thì. & Chủ động 100\% thời gian, biết chắc chắn lịch được nhận chỉ trong 39 giây. \\ \hline
Thất lạc dặn dò đặc biệt về tiền sử dị ứng xà phòng và tính cách. & Tự động đính kèm hồ sơ y tế \& Thẻ cảnh báo tương phản cao trên phiếu tiếp nhận. & Loại bỏ 100\% rủi ro quên dặn dò nguy hiểm, bảo vệ an toàn tuyệt đối cho thú cưng. \\ \hline
Mù mờ về tiến độ chăm sóc, lo âu suốt 2--3 tiếng, phải gọi điện hỏi. & Thanh tiến trình 4 mốc thời gian thực kết hợp thông báo đẩy khi chuyển giai đoạn. & Minh bạch thông tin, mang lại sự an tâm tuyệt đối cho chủ nuôi trong giờ làm việc. \\ \hline
Không có lịch sử lưu trữ, thông tin sản phẩm và hóa đơn bị phân tán. & Nhật ký chăm sóc cá nhân hóa, lưu hóa đơn điện tử \& Nút Rebook 1 chạm. & Dễ dàng tra cứu lại phác đồ chăm sóc cũ và tái đặt lịch định kỳ nhanh chóng. \\ \hline
\end{tabularx}
\caption{Tổng hợp các cải tiến thiết kế then chốt của hệ thống}
\label{tab:final_improvements}
\end{table}

\subsection{So sánh trước và sau cải tiến}
Đối chiếu định lượng và định tính giữa quy trình truyền thống (As-Is) và hệ thống đề xuất (To-Be):
\begin{table}[H]
\centering
\begin{tabularx}{\textwidth}{|l|X|X|}
\hline
\textbf{Tiêu chí so sánh} & \textbf{Quy trình truyền thống (As-Is)} & \textbf{Hệ thống mới đề xuất (To-Be)} \\ \hline
\textbf{Thời gian đặt lịch} & 15 phút -- vài tiếng (chờ tin nhắn phản hồi) & 39 giây (xác nhận tức thì trên ứng dụng) \\ \hline
\textbf{Xác suất sót dặn dò dị ứng} & Cao (phụ thuộc ghi nhớ và truyền miệng) & 0\% (tự động gắn cảnh báo từ hồ sơ số) \\ \hline
\textbf{Kênh cập nhật tiến độ} & Bị động (chủ nuôi phải gọi điện hỏi thăm) & Chủ động (thông báo đẩy 4 mốc thời gian thực) \\ \hline
\textbf{Khả năng tra cứu lịch sử} & Rời rạc, dễ thất lạc trong tin nhắn chat & Tập trung, minh bạch chi phí và phác đồ dịch vụ \\ \hline
\textbf{Mức độ hài lòng chung} & Thấp, thường xuyên lo lắng và ức chế & 4.16 / 5.0 (Tương đương 82.5 điểm SUS) \\ \hline
\end{tabularx}
\caption{So sánh hiệu quả trước và sau khi triển khai hệ thống}
\label{tab:before_after_comparison}
\end{table}
